\section{Renewal theory of the mean yield}\label{app:D.44}

The mean pair number of a cascade grows in proportion to the deposited energy, $\bar n(\omega) = \omega/\varepsilon_{\mathrm{eh}} + o(\omega)$ (\cref{eq:18.2-linear}). For a single chain of independent losses this is the key renewal theorem, \cref{eq:C.11-linear}. A cascade splits its energy at every ionization and has no common loss law, and we prove the leading term by following one unit of energy through the tree. The argument needs energy conservation and the Markov property of \cref{sec:18.1-branching}; it does not need the independence of siblings.

\subsection{The tagged-unit identity}

Start a cascade with one carrier of kinetic energy $E$ and let $N_{\mathrm{sub}}(E)$ be the number of ionizations in it. Represent the energy of the first carrier by the interval $[0, E]$. At every collision, divide the interval of the colliding carrier, of length $x$, into consecutive pieces. An ionization makes pieces of lengths $x_1$, $x_2$, $x_3$ for the three products and one of length $\Eg$ for the new pair. A phonon emission makes a piece of length $x - \omega_{\mathrm{AB}}$ for the carrier and one of length $\omega_{\mathrm{AB}}$ for the phonon. Energy conservation makes the pieces fill the interval exactly. A terminal carrier keeps its piece. Every point $u \in [0, E]$ then has a line: the sequence of carriers whose pieces contain $u$, ending when $u$ falls into a pair or phonon piece or its carrier becomes terminal. Let $x_1 = E, x_2, \dots$ be the kinetic energies of the carriers on the line of $u$ at its successive collisions, ionizations and phonon emissions alike; only the ionizations carry weight below.

Each ionization of a carrier with kinetic energy $x$ lies on the lines of exactly the $x$ units of its piece. Giving each of them the weight $1/x$ therefore counts the ionization once:
\begin{equation*}
  N_{\mathrm{sub}}(E) = \int_0^E\dd u\,S(u) , \qquad
  S(u) = \sum_{j\,\in\,\mathrm{ii}(u)}\frac{1}{x_j} ,
\end{equation*}
with $\mathrm{ii}(u)$ the collisions on the line of $u$ that are ionizations. Taking the mean and exchanging it with the integral,
\begin{equation}
  \langle N_{\mathrm{sub}}(E)\rangle = E\,\langle S\rangle_u ,
  \label{eq:D.44-tagged}
\end{equation}
with $u$ uniform on $[0, E]$ and the mean taken over $u$ and the cascade. This is \cref{eq:18.2-tagged}.

The line of a uniform unit is a Markov chain. Suppose that its carrier has kinetic energy $x$ and that, given everything before, $u$ is uniform on the carrier's piece. This holds for the first carrier. At the next collision the carrier's outcome is drawn from its kernel, independently of $u$ by the Markov property. The unit $u$ then lies in the piece of product $i$ with probability $x_i/x$, in the pair piece with probability $\Eg/x$, or in the phonon piece with probability $\omega_{\mathrm{AB}}/x$, and given that it lies in a piece it is uniform on it. The uniformity therefore propagates, and the line moves from $x$ to a product of energy $x_i$ with probability $x_i/x$ times the probability of that outcome, or leaves the chain. Write $\xi = \ln(x_{j+1}/x_j)$ for one step of $\ln x$ and $K_x$ for its law, with an extra state for leaving. Every step has $\xi \le 0$, because a product never carries more energy than its parent.

For a check, let a piece of length $x \ge 1$ split at a uniform point into two pieces, with no pair and no phonon, until every piece is shorter than 1. The mean number of splits $M(x)$ obeys $M(x) = 1 + (2/x)\int_1^xM(y)\,\dd y$ for $x \ge 1$, and $M(x) = 2x - 1$ solves it, since $(2/x)\int_1^x(2y - 1)\,\dd y = 2(x - 1)$. A unit lands in a piece of share $v$ with the size-biased density $2v$ on $[0, 1]$, so the walk of $\ln x$ has mean step $\int_0^12v\ln v\,\dd v = -\tfrac12$, and \cref{eq:D.44-tagged} gives $\langle S\rangle_u = M(x)/x = 2 - 1/x$, which tends to 2: one split per half unit of energy.

\subsection{Hypotheses}

Let $x_{\mathrm{cut}}$ be a cut energy and $\tau$ the first collision of the line at or below $x_{\mathrm{cut}}$. Four conditions are needed.
\begin{enumerate}
  \item \emph{Scale-free limit, approached fast enough.} There is a law $K_\infty$ of $\xi$, without leaving, such that $\|K_x - K_\infty\|_{\mathrm{TV}} \le C(\Eg/x)^\varkappa$ for some $C, \varkappa > 0$, with $\|\cdot\|_{\mathrm{TV}}$ the total-variation distance.
  \item \emph{No lattice.} $K_\infty$ is not concentrated on a lattice $\{0, -d, -2d, \dots\}$.
  \item \emph{Finite mean.} $K_\infty$ has a finite mean $-\mu_{\mathrm{ren}}$; by the second condition and $\xi \le 0$, $\mu_{\mathrm{ren}} > 0$.
  \item \emph{Regularity near threshold.} The mean weight collected at or below the cut, $h_{\mathrm{c}}(z) = \langle\sum_{j \ge \tau,\,j\in\mathrm{ii}}1/x_j\,|\,x_\tau = z\rangle$, is bounded uniformly in $x_{\mathrm{cut}}$ and continuous in $z$.
\end{enumerate}
In the random-$k$ cascade of \cref{sec:18.1-reductions} the shares of the products tend to a law independent of $x$: the Dirichlet law with parameters $3/2$ without a hole cutoff, and, with the hole density of states cut at $B_{\mathrm{h}}$, the Beta law with parameters $3/2$, $3/2$ for the two electrons, since then $\mathfrak{P}_{\mathrm{e}}(y) \propto y^2$ at $y \gg B_{\mathrm{h}}$ and the hole's share tends to zero. The probability of leaving into a pair is $\Eg/x$, and the phonon branch falls as $x^{-3}$ without the cutoff and as $x^{-3/2}$ with it; in both cases the leaving probability dominates and the first condition holds with $\varkappa = 1$. Both limit laws are continuous, which gives the second, and the mean of the logarithm of a share is finite for it, which gives the third. The fourth needs a finite mean number of collisions between $\Eg$ and $x_{\mathrm{cut}}$ and moments that are continuous in the energy of entry; both hold for the recursions \cref{eq:D.43-m1,eq:D.43-m2}, whose ionization probability vanishes continuously at threshold.

\subsection{Proof}

Split $S = S_> + S_\le$ into the ionizations above and at or below $x_{\mathrm{cut}}$.

\emph{Visits.} By the first and third conditions there are $\delta, p > 0$ and $x_1$ such that $P(\xi \le -\delta) \ge p$ under $K_x$ for every $x \ge x_1$. Since $\ln x_j$ never increases, the collisions the line spends in any interval of $\ln x$ of unit length above $\ln x_1$ form at most $1/\delta + 1$ runs, each of geometrically distributed length with parameter $p$. Their mean number is therefore at most a constant $C_{\mathrm{v}}$ independent of $E$. For $x_{\mathrm{cut}} \ge x_1$, summing over the intervals $[\ln x_{\mathrm{cut}} + k, \ln x_{\mathrm{cut}} + k + 1)$, in which each term of $S_>$ is at most $\ee^{-k}/x_{\mathrm{cut}}$,
\begin{equation*}
  \langle S_>\rangle \le C_{\mathrm{v}}\sum_{k \ge 0}\frac{\ee^{-k}}{x_{\mathrm{cut}}} = O\Bigl(\frac{1}{x_{\mathrm{cut}}}\Bigr) ,
\end{equation*}
uniformly in $E$.

\emph{Coupling.} Run the true line and a walk with steps drawn from $K_\infty$ from the same start, coupled at each step so that they take the same step with probability $1 - \|K_{x_j} - K_\infty\|_{\mathrm{TV}}$. The $K_\infty$ walk obeys the same visit bound. The probability that the two separate before reaching $x_{\mathrm{cut}}$ is at most the mean of $\sum_{x_j > x_{\mathrm{cut}}}\|K_{x_j} - K_\infty\|_{\mathrm{TV}}$, which by the first condition and the visit bound is $O((\Eg/x_{\mathrm{cut}})^\varkappa)$.

\emph{Undershoot.} Under $K_\infty$, $-\ln x_j$ is a renewal chain with independent, identically distributed, non-lattice increments of finite mean $\mu_{\mathrm{ren}}$. The undershoot $\ln x_{\mathrm{cut}} - \ln x_\tau$ is the overshoot of that chain over the level $\ln(E/x_{\mathrm{cut}})$. By the key renewal theorem, \cref{eq:C.11-key}, its law converges to the limit of \cref{eq:C.11-undershoot} as $E/x_{\mathrm{cut}}$ grows without bound. By the fourth condition $h_{\mathrm{c}}$ is bounded and continuous, so the mean of $h_{\mathrm{c}}(x_\tau)$ under $K_\infty$ converges.

\emph{Limit.} At $\tau$ the line restarts from $x_\tau$, so $\langle S_\le\rangle = \langle h_{\mathrm{c}}(x_\tau)\rangle$. For the true line this differs from the $K_\infty$ value by at most $2\sup h_{\mathrm{c}}$ times the separation probability. Together,
\begin{equation*}
  \limsup_{E\to\infty}\langle S\rangle_u - \liminf_{E\to\infty}\langle S\rangle_u \le O\Bigl(\frac{1}{x_{\mathrm{cut}}}\Bigr) + O\Bigl(\Bigl(\frac{\Eg}{x_{\mathrm{cut}}}\Bigr)^\varkappa\Bigr) ,
\end{equation*}
for every $x_{\mathrm{cut}}$. Letting $x_{\mathrm{cut}} \to \infty$, the limit of $\langle S\rangle_u$ exists. It is finite by the visit bound and positive whenever the cascade ionizes at all; call it $1/\varepsilon_{\mathrm{eh}}$. By \cref{eq:D.44-tagged}, $\langle N_{\mathrm{sub}}(E)\rangle = E/\varepsilon_{\mathrm{eh}} + o(E)$, and with the seed pair counted and $E = \omega - \Eg$ for a seed electron, $\bar n(\omega) = \omega/\varepsilon_{\mathrm{eh}} + o(\omega)$.

The proof claims no rate for the $o(\omega)$; the constant term is open (\cref{sec:18.2-open}). The lattice of energy losses in a one-phonon model does not affect the leading term, because the steps of $K_\infty$ come from continuous shares; it enters only the constant term, through the threshold region.
